% !TEX root = ../main.tex
\section{페이지랭크와 순위 상관의 기본 성질}
\label{ch:appendix-theory}

\ref{ch:literature}--\ref{ch:results}장은 여기에 정리한 기본 성질에 기대고 있어요. 이 내용은 널리 알려진 표준적인 것이에요. 페이지랭크(PageRank)의 성질은 Page 등\cend{15}과 Langville과 Meyer\cend{23}를 따르고, 순위 상관(두 줄 세우기가 얼마나 비슷한지 재는 수)은 Spearman\cend{20}과 Kendall\cend{21}을 따라요.

\dissertationSubheading[sec:pr-existence]{페이지랭크의 존재와 유일성}

$P$를 점 $|V|$개 위의 행 확률 행렬(각 행의 합이 1인 행렬)이라고 해요. 필요하면 식~\eqref{eq:transition}의 방법으로 매달린 점(나가는 화살표가 없는 점)을 고친 다음의 행렬이에요. $d\in(0,1)$일 때 식~\eqref{eq:google-matrix}의 구글 행렬은 다음과 같아요.
\[
G_{\mathrm{PR}}=d P^{\top}+\frac{1-d}{|V|}\mathbf{1}\mathbf{1}^{\top},
\]
이 행렬은 모든 칸이 양수이고, 각 열의 합이 1인 열 확률 행렬이에요. 양수 행렬에 대한 페론--프로베니우스(Perron--Frobenius) 정리에 따르면, $G_{\mathrm{PR}}$에는 크기가 1인 고윳값(행렬이 벡터를 몇 배로 늘리는지 나타내는 수)이 하나뿐이에요. 그리고 그에 맞는 고유벡터 $\mathbf{r}$는 모든 칸이 0보다 커요. 우리는 이 벡터를 $\sum_i r_i=1$이 되도록 크기를 맞출 수 있어요. 따라서 감쇠 계수(화살표를 따라갈 확률)와 매달린 점 규칙을 정하고 나면, 방향이 있는 그래프는 모두 단 하나의 페이지랭크 점수를 가져요\cend{23}.

이 논문의 재시작 벡터 가운데 두 개는 고르지 않아요. AWP(적응형 가중 페이지랭크)의 재시작 벡터는 활발함(활동 정도)에 비례하므로, 아무것도 보내지 않은 주소에서는 0이에요. S-PR(보증 씨앗 페이지랭크)의 재시작 벡터는 허락 걷기의 점수를 합이 1이 되게 맞춘 것이므로, 허락 점수가 없는 송금 층 지갑에서는 0이에요. 그래서 두 구글 행렬 모두 모든 칸이 양수이지는 않아요. 그래도 고정점(계산을 되풀이해도 더는 바뀌지 않는 값)은 여전히 하나뿐이에요. 두 확률 벡터의 차이를 보면 재시작 항이 서로 지워지고, 연산자는 두 벡터 사이의 $\|\cdot\|_1$ 거리를 많아야 $d$배로 만들기 때문이에요. 하지만 일반적으로 이 고정점의 모든 칸이 0보다 크지는 않아요. 재시작 몫이 없고, 재시작 지갑에서 출발한 어떤 길도 닿지 않는 지갑은 점수가 계속 0으로 남기 때문이에요.

답이 하나뿐이라는 성질(유일성) 덕분에 엔도스랭크(EndorseRank)와 AWP를 분명하게 비교할 수 있어요. 매개변수를 정하고 나면, 두 방법은 저마다 주어진 화살표 묶음에 꼭 하나의 점수 벡터를 줘요. 그래서 $\mathbf{s}_{\mathrm{ER}}$와 $\mathbf{s}_{\mathrm{AWP}}$의 차이는 모두 두 방법의 화살표와 재시작 분포에서 생길 수밖에 없어요. 여러 풀이 프로그램 탓으로 돌릴 수는 없어요.

\dissertationSubheading[sec:pr-continuity]{화살표 무게에 대한 연속성}

$d\in(0,1)$를 고정하면, 페이지랭크는 $P$의 칸 값을 따라 연속적으로(끊김 없이) 바뀌어요. $\tilde P$가 같은 점들 위의 또 다른 행 확률 행렬이고 그 페이지랭크 벡터가 $\tilde{\mathbf{r}}$라면, 다음 부등식이 맞아요.
\[
\|\mathbf{r}-\tilde{\mathbf{r}}\|_1\le\frac{d}{1-d}\,\max_i\bigl\|P_{i\cdot}-\tilde P_{i\cdot}\bigr\|_1,
\]
이 한계는 $|V|$에 달려 있지 않아요\cend{23}.

$P$에 대한 연속성은 화살표 무게에 대한 연속성과 같지 않아요. 행 $P_{i\cdot}=w_{i\cdot}/o_i$는 나가는 무게의 합 $o_i$가 0보다 큰 동안에만 무게를 따라 연속적으로 움직여요. 어떤 점이 처음으로 나가는 화살표를 얻거나 마지막 나가는 화살표를 잃으면, 그 행은 고른 매달린 행과 나가는 이웃에게 몰린 행 사이를 훌쩍 건너뛰어요. 관련된 무게가 아무리 작아도 그래요. \ref{sec:pagerank-worked-example}절의 계산 예에서, 무게가 $10^{-12}$인 화살표 $D\to A$를 더하면 $\mathbf{r}$가 $\|\cdot\|_1$ 노름(칸마다 차이의 크기를 모두 더한 값)으로 $0.37$만큼 움직이고, $r_A$는 $0.1375$에서 $0.3202$로 올라가요.

감쇠 계수를 바꿔 보는 실행은 연속성으로 설명돼요. 이 실행은 $P$를 그대로 두고 $d$만 매끄럽게 바꾸는데, $d<1$이면 페이지랭크 벡터는 $d$를 따라 연속적으로 바뀌기 때문이에요. 상위 토큰(가장 많이 쓰인 토큰)만 남기는 거르기는 연속성으로 설명되지 않아요. 이 거르기는 화살표를 지우는데, 그러면 어떤 지갑은 나가는 화살표가 하나도 없게 될 수 있어요. 그래서 이 점검이 얼마나 안정적인지는 오로지 자료에 달려 있어요(표~\ref{tab:robustness-tokens}). 실제 자료에서 엔도스랭크의 허락 지표 가족(비슷한 지표 묶음) 평균은 $d\in\{0.75,0.85,0.95\}$에서 소수점 아래 세 자리까지 바뀌지 않았어요(표~\ref{tab:robustness-damping}). 코호트의 대부분이 동점이니, 순위는 거의 움직일 수 없었을 거예요.

\dissertationSubheading[sec:pr-sparsity]{희소성과 반복 한 번의 비용}

0이 아닌 화살표가 $m=|E|$개 있으면, 희소한(듬성듬성한) 곱 $P^{\top}\mathbf{x}$에는 $O(m)$번의 셈이 들어요. 화살표 목록으로 희소 연산자를 만드는 것도 $O(m)$이에요. 그래서 $I$번 하는 거듭제곱 반복(같은 계산을 되풀이하기)은 $O(m+mI)$예요. 매칭 코호트(모든 점수가 똑같이 점수를 매기는 지갑 묶음)에서 송금 그래프는 보증 그래프보다 화살표가 약 $9.3$배 많아요. 그래서 $I$가 똑같더라도 송금 방법은 대략 한 자릿수(약 열 배)만큼 더 오래 걸릴 거예요. 실제로 잰 호출에서는 연산자를 조립하는 데 시간이 대부분 들고, 반복에는 조금만 들어요(\ref{sec:benchmark-results}절). 그래서 잰 시간의 비율 $9.2$는 화살표 수의 비율을 따라가요.

\dissertationSubheading[sec:spearman-properties]{스피어먼의 \texorpdfstring{$\rho$}{rho}의 성질}

스피어먼의 로($\rho$)는 순위로 계산한 피어슨 상관(두 값이 함께 커지고 작아지는 정도)이에요. 이 값은 $[-1,1]$ 안에 있어요. 그리고 두 변수 가운데 어느 쪽을 늘 커지기만 하는 방식으로 바꾸어도 달라지지 않아요. $\rho=1$이면 순서가 완전히 같은 방향으로 맞고, $\rho=-1$이면 완전히 거꾸로 맞는다는 뜻이에요. 같은 값(동점)에는 중간 순위(동점끼리 나눠 갖는 평균 순위)를 줘요. 평가 처리 과정은 동점을 제대로 다루는 표준 구현을 써요. 평판 점수와 대리 지표(대신 재는 값)는 연속적인 값이고 꼬리가 두꺼워서(아주 큰 값이 드물게 나와서), 원래 값에 대한 피어슨 상관보다 순위에 바탕을 둔 관계 재기가 더 잘 맞아요. AWP 기준선(비교 기준) 논문도 마찬가지로 자기 순위를 순위 통계로 평가해요\cend{3}.

\dissertationSubheading[sec:kendall-properties]{켄달의 \texorpdfstring{$\tau$}{tau}의 성질}

켄달의 타우($\tau$)는 순서가 맞는 쌍의 수 $n_c$와 순서가 어긋나는 쌍의 수 $n_d$를 견주어요. 동점 보정을 하지 않으면, $n_0=n(n-1)/2$일 때 $\tau_a=(n_c-n_d)/n_0$는 무작위로 뽑은 두 지갑을 두 순위가 같은 방향으로 줄 세울 확률에서 반대 방향으로 줄 세울 확률을 뺀 값이에요\cend{21}. 보고한 $\tau$는 모두 식~\eqref{eq:kendall}의 동점 보정 $\tau_b$예요. 동점이 없으면 두 값은 같지만, 여기서는 동점이 아주 많아요. 엔도스랭크는 매칭 코호트에서 서로 다른 값을 $99$개만 가져요. 그래서 모든 지갑 쌍의 $67.3\%$가 엔도스랭크에서 동점인데, AWP에서는 $0.5\%$뿐이에요. 그리고 엔도스랭크와 AWP 사이에 보고한 $\tau_b=0.360$은 $\tau_a=0.206$과 짝이 돼요. 동점은 $\tau_b$의 위쪽 한계도 정해요. 엔도스랭크가 얻을 수 있는 가장 큰 값은 허락 대리 지표에 대해서는 약 $0.38$이에요. 허락 대리 지표는 엔도스랭크가 떼어 내는 바로 그 131개 지갑에서만 0보다 커요. 동점이 없는 대리 지표에 대해서는 약 $0.57$이에요. 그래서 허락 가족에서의 $0.375$(표~\ref{tab:alignment-hybrid})는 두 순위가 떼어 내는 쌍에서는 거의 완전히 맞는다는 뜻이에요. 그리고 이 값은 천장(얻을 수 있는 가장 큰 값)이 다른 계수와 곧바로 비교할 수 없어요.

$n$이 크면 $\tau$와 $\rho$는 보통 부호가 같지만, 크기는 달라요. 엔도스랭크나 AWP를 대리 지표와 짝지은 것 가운데 $|\tau|>0.02$인 짝들에서, 비율 $\rho/\tau$는 $1.01$과 $1.49$ 사이에 있고 늘 $\rho$가 더 커요. AWP 기준선이 그렇게 하듯이, 우리도 둘 다 보고해요\cend{3}. 우리의 주된 요약 통계는 짝지은 부트스트랩(다시 뽑기) 구간을 붙인 가족 평균 $\tau$예요\cend{34}.

\dissertationSubheading[sec:construct-validity-note]{구성 타당도에 대한 메모}

Cronbach와 Meehl\cend{22}은 기준 관련 타당도(정답 기준과 얼마나 맞는지)와 구성 타당도(점수가 재려는 것을 정말 재는지)를 나누어요. 확실한 정답 기준이 되는 채무 불이행(빌린 돈을 갚지 못함) 라벨이 없으므로, 이 연구는 구성개념(점수가 재려고 하는 생각)을 중심에 둔 방법을 써요. 각 점수는 그 점수의 개념 정의를 담은 검증 점검으로 시험해요(예를 들어 엔도스랭크는 허락을 받는 것, AWP는 들어오는 송금을 받는 것). 그리고 두 라벨(나중에 채점에 쓰는 정답 값), 곧 새 허락 쌍과 새 송금자가 있는 홀드아웃(떼어 두었다가 나중에 맞혀 보는) 기간으로도 시험해요. 같은 구성개념 안에서 $\tau$가 높으면 이 해석을 뒷받침해요. 반대로 서로 다른 구성개념 사이에서 $\tau$가 낮다고 해서 저절로 실패를 뜻하지는 않아요. 그것은 판별 타당도(서로 다른 것을 다르게 재는지)를 보여 주는 것일 수도 있어요.

\dissertationSubheading[sec:numerical-summary-sheet]{주요 값을 보고한 곳}

아래 목록은 주요 값마다 그 값을 내놓은 표를 알려 줘요.

\begin{center}
\fitwidth{%
\begin{tabular}{ll}
\toprule
값 & 보고한 곳 \\
\midrule
매칭 코호트 크기, $|V|$, $|E|$, 실행 시간, 표준편차, 메모리, 반복 수(방법 네 개) & 표~\ref{tab:benchmark-runtime} \\
코호트 안 규모 키우기 단계 & 표~\ref{tab:benchmark-scaling-incohort} \\
풀 부분 표본 단계($N=99{,}080$) & 표~\ref{tab:benchmark-scaling} \\
95\% 신뢰구간을 붙인 가족 평균 $\tau$, 모든 점수 & 표~\ref{tab:alignment-hybrid} \\
C-PR과 S-PR의 짝지은 $\Delta\tau$ 대비 & 표~\ref{tab:tau-diff-hybrid} \\
방법 사이 $\tau$; 그래프 밖 지갑을 외톨이 점으로 둔 경우 & 표~\ref{tab:tie-sensitivity} \\
감쇠 계수, 상위 토큰, 표본 정의 바꿔 보기 & 표~\ref{tab:robustness-damping}, \ref{tab:robustness-tokens}, \ref{tab:robustness-sample-size} \\
스펜더 홀드아웃($n=1{,}335$) & 표~\ref{tab:holdout-spenders}, \ref{tab:holdout-diff}과~\ref{tab:holdout-receiving} \\
탐색적 트레이더 실행 & 표~\ref{tab:holdout-traders} \\
등록한 6월부터 8월까지의 재현 & \ref{sub:fresh-results}절 \\
EndorseRank와 AWP 비교, 대리 지표별 결과 & 부록~\ref{ch:appendix-er-awp} \\
관찰 기간 & 2025-12-01 -- 2026-05-31(표~\ref{tab:parameter-summary}) \\
\bottomrule
\end{tabular}
}
\end{center}

여기 적은 값은 모두 평가 요약에도 있고, \ref{ch:results}장이나 부록의 표에도 나와요.
